\clearpage
\VerdierPhysicalPageAnchor{0049}%
各 $\sigma\in G$ に対して、図式

\VerdierSemanticAnchor{ch1-1-4-5-eq4}
\begin{center}
  \resizebox{0.78\textwidth}{!}{%
    $\begin{tikzcd}[ampersand replacement=\&,column sep=6em,row sep=4em]
      F_1T(\sigma)
        \arrow[r,"{p_1(\sigma)}"]
        \arrow[d,"{u*T(\sigma)}"']
        \& T'(\sigma)F_1
          \arrow[d,"{T'(\sigma)*u}"] \\
      F_2T(\sigma)
        \arrow[r,"{p_2(\sigma)}"]
        \& T'(\sigma)F_2
    \end{tikzcd}$%
  }
  \hfill\textup{(1.4.5.4)}
\end{center}
は $\beta(\sigma)$-可換である。

\VerdierSemanticAnchor{ch1-1-4-6}
\noindent\textbf{定義 1.4.6.}\quad
$(\mathcal C,G,T,c)$ を $G$-圏、$(\mathcal C',G,T',c')$ を加法的な
$G$-圏とする。$\alpha_1,\alpha_2$ を
$Z^2(G,\mathbb Z/2\mathbb Z)$ の正規化された二つの非斉次余サイクルとし、
$\beta\colon G\to\mathbb Z/2\mathbb Z$ を

\VerdierSemanticAnchor{ch1-1-4-6-eq1}
\begin{equation}
  \alpha_2(\sigma_1,\sigma_2)\alpha_1^{-1}(\sigma_1,\sigma_2)
  =\beta(\sigma_2)\beta(\sigma_1\sigma_2)^{-1}\beta(\sigma_1)
  \tag{1.4.6.1}
\end{equation}
を満たす写像とする。

\VerdierSemanticAnchor{ch1-1-4-6-beta-morphism}
最後に
$(F_1,p_1),(F_2,p_2)\colon(\mathcal C,G,T,c)
\to(\mathcal C',G,T',c')$
を、それぞれ余サイクル $\alpha_1,\alpha_2$ で捩られた二つの $G$-関手と
する。関手の射 $u\colon F_1\to F_2$ で、各 $\sigma\in G$ に対して図式
(1.4.5.4) が $\beta(\sigma)$-可換となるものを、\emph{関手 $(F_1,p_1)$ から
関手 $(F_2,p_2)$ への $\beta$-射}という。

\VerdierSemanticAnchor{ch1-1-4-7}
\noindent\textbf{1.4.7.}\quad
(1.4.5) で示したように、正規化された二つの切断を選ぶことにより、符号群の
標準作用を誘導する $H$-関手の射に、対応する捩れ $G$-関手の間の
$\beta$-射を対応させられる。この対応が全単射であることは直ちに確かめられる。
さらに、符号群の標準作用を誘導する三つの $H$-関手
$(F_1,m_1)$、$(F_2,m_2)$、$(F_3,m_3)$、正規化された三つの切断
$s_1,s_2,s_3$、二つの $H$-関手の射 $u,v$ とその合成 $vu$：

\VerdierSemanticAnchor{ch1-1-4-7-eq1}
\begin{equation}
  \begin{aligned}
    u&\colon(F_1,m_1)\longrightarrow(F_2,m_2),\\
    v&\colon(F_2,m_2)\longrightarrow(F_3,m_3),\\
    vu&\colon(F_1,m_1)\longrightarrow(F_3,m_3)
  \end{aligned}
  \tag{1.4.7.1}
\end{equation}
が与えられたとき、射

\VerdierSemanticAnchor{ch1-1-4-7-eq2}
\begin{equation}
  \begin{aligned}
    u&\colon(F_1,p_1)\longrightarrow(F_2,p_2),\\
    v&\colon(F_2,p_2)\longrightarrow(F_3,p_3),\\
    vu&\colon(F_1,p_1)\longrightarrow(F_3,p_3)
  \end{aligned}
  \tag{1.4.7.2}
\end{equation}

\VerdierSemanticAnchor{ch1-1-4-7-open-seam}
